% Adaptación de «Balances térmicos, renovación y radiación» del HTML.
\begin{table}[p]
\centering
\caption{Nomenclatura y signos de los balances del recinto sin carga.}
\label{met:tabla_pinn_balances}
\begingroup
\small
\setlength{\tabcolsep}{4pt}
\renewcommand{\arraystretch}{1.12}
\begin{tabularx}{\textwidth}{@{}>{\raggedright\arraybackslash}p{0.23\textwidth}>{\raggedright\arraybackslash}X>{\raggedright\arraybackslash}p{0.21\textwidth}@{}}
\toprule
Símbolo & Significado y convención de signo & Unidad \\
\midrule
$C_c,C_b,C_a(W)$ & Capacidades térmicas de cubierta, base y aire húmedo. & J/K \\
$G_{ca},G_{ba}$ & Conductancias de intercambio cubierta--aire y base--aire; pueden depender del modo de ventilación. & W/K \\
$Q_{\mathrm{sol},c},Q_{\mathrm{sol},b}$ & Potencias solares absorbidas; positivas hacia el nodo correspondiente. & W \\
$Q_{c,\mathrm{ext}},Q_{b,\mathrm{ext}}$ & Pérdidas exteriores; positivas al salir de cubierta o base. & W \\
$Q_{b\to c}^{\mathrm{rad}}$ & Intercambio radiativo neto; positivo desde la base hacia la cubierta. & W \\
$F_{\mathrm{ext}}$ & Flujo exterior de aire seco con entrada y salida equilibradas en la reducción de inventario constante. & kg de aire seco/s \\
$m_{da}$ & Inventario de aire seco, constante en las Ecuaciones~\ref{met:red_cubierta}--\ref{met:red_humedad}. & kg de aire seco \\
$c_{da},c_{pv}$ & Calores específicos del aire seco y del vapor. & J/(kg\,K) \\
$Q_{\mathrm{fan}}$ & Fracción de la potencia del equipo transferida al dominio modelado. & W \\
$A_b,A_c$ & Áreas de base y cubierta utilizadas en el cierre radiativo. & m$^2$ \\
$\epsilon_b,\epsilon_c;\ F_{bc}$ & Emisividades y factor de forma desde la base hacia la cubierta. & Adimensional \\
$\sigma$ & Constante de Stefan--Boltzmann. & W/(m$^2$\,K$^4$) \\
$\dot{(\,)};\ \partial/\partial s$ & Derivada temporal del balance y derivada de la red respecto del tiempo local, con condiciones del segmento fijas. & Estado por segundo \\
\bottomrule
\end{tabularx}
\endgroup
\par\smallskip
\begin{minipage}{\textwidth}\footnotesize
La irradiancia $G$ se expresa en W/m$^2$, mientras $G_{ca}$ y $G_{ba}$ son conductancias en W/K. El flujo másico $F_{\mathrm{ext}}$ y el factor de forma $F_{bc}$ representan magnitudes distintas. Los subíndices $c,b,a$ designan cubierta, base y aire; $\mathrm{amb}$ y $\mathrm{ext}$ designan ambiente e intercambio exterior.
\end{minipage}
\end{table}
